% notes/v5/lit_framing.tex -- literature framing for the Applied Mathematics Letters version.
% Draft, 2026-09-28. \cite keys refer to notes/v5/refs_letter.bib. Plain LaTeX; no macros from
% paper/main.tex are needed. Notation: h_\Omega bulk mesh size (Scott's penalty is \mu/h_\Omega),
% h_\Gamma longest chord, \Gamma_h the inscribed polygon, \bar p the mean of p over \Gamma.
% Paragraphs A1-A2 = history, B = positioning, C = novelty, D = scope. Target length: about one page.

% ---------------------------------------------------------------- (A) history and background
\paragraph{Exactly divergence-free elements.}
Scott and Vogelius~\cite{ScottVogelius1985} showed that continuous piecewise-$P_k$ velocities with discontinuous piecewise-$P_{k-1}$ pressures satisfy $\operatorname{div}V_h=\Pi_h$ for $k\ge4$ on meshes without singular vertices; the inf-sup constant was later made explicit in terms of the local mesh geometry~\cite{GuzmanScott2019}. The discrete velocity is then divergence-free pointwise, not only weakly. This gives exact mass conservation. On fitted meshes with strongly imposed no-slip data it also gives pressure robustness: the velocity error does not depend on the pressure, which matters when pressure gradients are large or the viscosity is small~\cite{JohnEtAl2017}. Split meshes give similar pairs of lower degree~\cite{GuzmanNeilan2018}; see~\cite{Neilan2020} for a review.

\paragraph{Curved walls and Nitsche's method.}
Exact incompressibility has a cost at a curved wall. A conforming divergence-free field that vanishes on a polygon has a vanishing gradient at every boundary vertex~\cite{ScottPrize,GjerdeScott2024}. If the wall is replaced by an inscribed polygon and no-slip is imposed strongly, the computed gradient is therefore off by $O(1)$ at each vertex. The $H^1$ error of Scott--Vogelius elements then decays only like $h^{1/2}$~\cite{ScottPrize,GjerdeScott2024}, compared with $h^{3/2}$ for scalar problems~\cite{BergerScottStrang1972,Scott1975}. One remedy is to fit the geometry, with isoparametric or Piola-mapped elements~\cite{NeilanOtus2021,DurstNeilan2024}. Gjerde and Scott~\cite{GjerdeScott2024} keep the polygon and impose the wall condition weakly by Nitsche's method~\cite{Nitsche1971}, with penalty $\mu/h_\Omega$ and no pressure term, and report errors orders of magnitude below those of strong imposition. Scott observed $H^1$ errors of order $h^{3/2}$. He offered a prize for a published proof that Scott--Vogelius--Nitsche converges, with an expected error $h_\Gamma^{3/2}+h_\Omega^k$, and asked: ``If this does not hold, why not?''~\cite{ScottPrize}.

% ---------------------------------------------------------------- (B) positioning
\paragraph{Related work.}
For Stokes, the consistent Nitsche form uses the full traction $\sigma(u,p)n$. It therefore contains the pressure term $\langle p,v\cdot n\rangle$~\cite{FreundStenberg1995,Stenberg1995,Becker2002,BurmanHansbo2014}, and this term is recognised as essential for consistency~\cite{BurmanHansboLarson2019Stokes}. The Gjerde--Scott form omits it. On the geometric side, boundary-value correction~\cite{BrambleDupontThomee1972} has been combined with Nitsche's method for cut elements~\cite{BurmanHansboLarson2018,BurmanHansboLarson2019Stokes} and for polygonal approximations~\cite{DupontGuzmanScott2025}. The shifted boundary method~\cite{MainScovazzi2018} is a related approach. None of these works uses exactly divergence-free spaces.
With such spaces, the placement of the pressure term decides whether exact incompressibility survives. If the term also appears in the continuity equation, $\operatorname{div}u_h$ vanishes only away from an $O(h)$ boundary strip~\cite{LiuNeilanOlshanskii2023}. If it appears only in the momentum equation, the constant pressure mode must be handled separately. Frachon, Nilsson and Zahedi~\cite{FrachonNilssonZahedi2024} show that fixing the pressure mean by a Lagrange multiplier makes $\operatorname{div}u_h$ equal to that multiplier. They avoid this by testing the momentum equation only with velocities of zero net flux. The work closest to ours is that of Liu, Neilan and Otus~\cite{LiuNeilanOtus2023}, for Scott--Vogelius elements on Clough--Tocher splits of a mesh inside a curved domain. They add a boundary multiplier that approximates the wall pressure up to a constant, together with a zero-flux constraint and a Taylor boundary correction. Their velocity is exactly divergence-free and converges at the optimal order $h^k$; without the constraint, they note, the method is ill posed. Related multiplier formulations for cut elements are given in~\cite{BurmanHansboLarson2024}. Exactly divergence-free methods also require discrete Dirichlet data with zero net flux~\cite{EickmannScottTscherpel2025}.

% ---------------------------------------------------------------- (C) novelty
\paragraph{Contribution.}
Our subject is the method as Gjerde and Scott printed it: no extra unknowns, no boundary correction, and the pressure itself on the wall. (i)~Omitting $\langle p,v\cdot n\rangle$ leaves the wall pressure to be balanced by the penalty rather than by a traction. The discrete flow then leaks through the polygon with normal velocity $(h_\Omega/\mu)(p-\bar p)$. We prove two-sided bounds: the slip error and the energy error, normalised by $(h_\Omega/\mu)\|p-\bar p\|_{L^2(\Gamma)}$ and $(h_\Omega/\mu)^{1/2}\|p-\bar p\|_{L^2(\Gamma)}$ respectively, both tend to $1$, and the $H^1$ rate is $1$. So Scott's expected estimate holds for this method if and only if the pressure is constant along the wall. (ii)~Adding the missing term makes the square Scott--Vogelius system singular, with the constant pressure as kernel. (iii)~Subtracting the wall mean, $\langle p_h-\bar p_{\Gamma_h}(p_h),v\cdot n\rangle$, gives a square system that is well posed for large enough penalty. It keeps $u_h$ exactly divergence-free and attains $h_\Gamma^{3/2}+h_\Omega^k$ for general Stokes data and every $k\ge4$; the $h_\Gamma^{3/2}$ term is sharp. (iv)~The penalty threshold must scale like $h_\Omega/\min_{e}|e|$. We do not claim the device in (iii) as new. Its velocity coincides with that of the zero-net-flux formulation of~\cite{FrachonNilssonZahedi2024}, applied to our setting; the two differ only in how the pressure constant is fixed. As far as we know, what is new is the analysis: the sharp leak characterisation (i), and the proof of the conjectured rate, with its sharpness, for Scott--Vogelius elements on an inscribed polygon.

% ---------------------------------------------------------------- (D) scope and variants
\paragraph{Scope.}
We analyse the symmetric Nitsche form of~\cite{GjerdeScott2024}. The leak comes from the missing pressure term, not from the symmetry of the viscous terms, so we expect it in nonsymmetric variants as well; we have not analysed these, nor penalty-free variants~\cite{FreundStenberg1995,Burman2012,BoiveauBurman2016}. The mechanism is the classical consistency error of boundary penalty methods~\cite{Babuska1973}. Here it acts only on the pressure, with vanishing penalty parameter $h_\Omega/\mu$, which is why the method converges at all. Nitsche's method with Robin-type data is treated in~\cite{JuntunenStenberg2009}.
Multiplier and boundary-correction methods~\cite{LiuNeilanOtus2023,BurmanHansboLarson2024,BrambleDupontThomee1972,DupontGuzmanScott2025} reach higher order than $h_\Gamma^{3/2}$, at the price of extra unknowns or Taylor corrections.
A common alternative restricts the pressure to $L^2_0$ and adds the plain term $\langle p_h,v\cdot n\rangle$. The continuity equation is then tested only with mean-zero pressures. Because $\operatorname{div}V_h=\Pi_h$, it gives only $\operatorname{div}u_h\equiv c$ with $c=|\Omega_h|^{-1}\int_{\partial\Omega_h}u_h\cdot n$ (compare~\cite[(5.9)]{FrachonNilssonZahedi2024}). The constant vanishes exactly when the discrete net flux through $\partial\Omega_h$ vanishes. Equivalently, it vanishes when this solution coincides with that of (iii), which happens exactly when the pressure of (iii) has zero mean on $\Gamma_h$. With weak imposition nothing forces the flux through $\Gamma_h$ to vanish, even for flux-corrected data. In our computations $c\neq0$.
Our proofs use $k\ge4$ Scott--Vogelius elements, and $k\ge2$ on Clough--Tocher splits, in two dimensions. Other divergence-free families, and cut or unfitted meshes, are not treated.
